\ifdefined\gkver\else\def\gkver{student}\fi
\documentclass[\gkver]{gaokaozhenti}
\begin{document}

\chapter{2012年浙江理综}
%% paperType: 理综物理部分

\begin{enumerate}
\item
%% number: 14
%% paperName: 2012年浙江理综第 14 题 6 分
%% typeId: 1
%% type: 单选题
%% chapter: 力物体的平衡
%% point: 多力平衡
%% method: 无
%% score: 6
%% degree: 650
%% duplicateId: 0
%% body:
如图所示，与水平面夹角为$30^{\circ}$的固定斜面上有一质量$m=1.0\Ukg$的物体。细绳的一端与物体相连，另一端经摩擦不计的定滑轮与固定的弹簧秤相连。物体静止在斜面上，弹簧秤的示数为$4.9\UN$。关于物体受力的判断（取$g=9.8\Umsq$），下列说法正确的是\xzanswer{A}
\onepicture[h=2.4cm, align=right]{figs/14.png}
\fourchoices[answer=A]
{斜面对物体的摩擦力大小为零}
{斜面对物体的摩擦力大小为$4.9\UN$，方向竖直向上}
{斜面对物体的支持力大小为$4.9\sqrt 3 \UN$，方向竖直向上}
{斜面对物体的支持力大小为$4.9\UN$，方向垂直斜面向上}
%% answer: A
%% memo:
\memoanswer{
由弹簧秤的示数为 $4.9\UN$ 可知细绳对物体的拉力大小也为 $4.9\UN$，刚好等于物体的重力沿斜面向下的分力，因此物体在重力、绳子的拉力和斜面的支持力下能够保持平衡状态，故选项 A 正确、B 错误；由平衡条件，斜面对物体的支持力 $F_N=mg\cos30^{\circ}=4.9\sqrt{3}\UN$，方向垂直斜面向上，故选项 C、D 错误。
}

\item
%% number: 15
%% paperName: 2012年浙江理综第 15 题 6 分
%% typeId: 1
%% type: 单选题
%% chapter: 曲线运动
%% point: 天体运动
%% method: 无
%% score: 6
%% degree: 650
%% duplicateId: 0
%% body:
如图所示，在火星与木星轨道之间有一小行星带。假设该带中的小行星只受到太阳的引力，并绕太阳做匀速圆周运动。下列说法正确的是\xzanswer{C}
\onepicture[h=2.4cm, align=right]{figs/15.png}
\fourchoices[answer=C]
{太阳对各小行星的引力相同}
{各小行星绕太阳运动的周期均小于一年}
{小行星带内侧小行星的向心加速度值大于外侧小行星的向心加速度值}
{小行星带内各小行星圆周运动的线速度值大于地球公转的线速度值}
%% answer: C
%% memo:
\memoanswer{
各小行星质量以及其与太阳的距离不一定相同，故太阳对各小行星的引力不一定相同，选项 A 错误；由 $a_n=\frac{GM}{r^2}$ 可知，内侧小行星与太阳的距离较近，向心加速度较大，选项 C 正确；各小行星绕太阳运动的半径均大于地球绕太阳运动的半径，由开普勒第三定律可知，各小行星绕太阳运动的周期均大于地球绕太阳运动的周期，选项 B 错误；根据 $G\frac{Mm}{r^2}=\frac{mv^2}{r}$，得 $v=\sqrt{\frac{GM}{r}}$，可知各小行星运动的线速度均小于地球公转的线速度，故选项 D 错误。
}

\item
%% number: 16
%% paperName: 2012年浙江理综第 16 题 6 分
%% typeId: 1
%% type: 单选题
%% chapter: 机械波
%% point: 波的图象
%% method: 无
%% score: 6
%% degree: 650
%% duplicateId: 0
%% body:
用手握住较长软绳的一端连续上下抖动，形成一列简谐横波。某一时刻的波形如图所示，绳上a、b两质点均处于波峰位置。下列说法正确的是\xzanswer{D}
\onepicture[h=2.8cm, align=right]{figs/16.png}
\fourchoices[answer=D]
{a、b两点之间的距离为半个波长}
{a、b两点振动开始时刻相差半个周期}
{b点完成全振动次数比a点多一次}
{b点完成全振动次数比a点少一次}
%% answer: D
%% memo:
\memoanswer{
由图可知，a、b 两点之间的距离为一个波长，振动开始时刻相差一个周期，a 点离波源近，完成全振动的次数比 b 点多一次，故选项 A、B、C 错误，选项 D 正确。
}

\item
%% number: 17
%% paperName: 2012年浙江理综第 17 题 6 分
%% typeId: 1
%% type: 单选题
%% chapter: 电路
%% point: 电功电功率
%% method: 近似与估算
%% score: 6
%% degree: 650
%% duplicateId: 0
%% body:
功率为$10\UW$的发光二极管（LED灯）的亮度与功率为$60\UW$的白炽灯相当。根据国家节能战略，2016年前普通白炽灯应被淘汰。假设每户家庭有2只$60\UW$的白炽灯，均用$10\UW$的LED灯替代，估算出全国一年节省的电能最接近\xzanswer{B}
\fourchoices[answer=B]
{$8\times10^{8}\UkWh$}
{$8\times10^{10}\UkWh$}
{$8\times10^{11}\UkWh$}
{$8\times10^{13}\UkWh$}
%% answer: B
%% memo:
\memoanswer{
每户两只白炽灯被两只 LED 灯替换后，节约电功率为 $120\UW-20\UW=100\UW$，每天按照五个小时用电计算，节电 $0.5\UkWh$，全年按照 360 天计算，节电 $180\UkWh$，全国按照四亿家庭计算，节电 $7.2\times10^{10}\UkWh$，故选项 B 正确。
}

\item
%% number: 18
%% paperName: 2012年浙江理综第 18 题 6 分
%% typeId: 2
%% type: 多选题
%% chapter: 功和能
%% point: 机械能守恒定律
%% method: 无
%% score: 6
%% degree: 450
%% duplicateId: 0
%% body:
由光滑细管组成的轨道如图所示，其中$AB$段和$BC$段是半径为$R$的四分之一圆弧，轨道固定在竖直平面内。一质量为$m$的小球，从距离水平地面高为$H$的管口$D$处静止释放，最后能够从$A$端水平抛出落到地面上，下列说法正确的是\xzanswer{BC}
\onepicture[h=3.0cm, align=right]{figs/18.png}
\fourchoices[answer=BC]
{小球落到地面时相对于$A$点的水平位移值为$2\sqrt {RH - 2{R^2}} $}
{小球落到地面时相对于$A$点的水平位移值为$2\sqrt {2RH - 4{R^2}} $}
{小球能从细管$A$端水平抛出的条件是$H>2R$}
{小球能从细管$A$端水平抛出的最小高度$H_{min}=\frac{5}{2}R$}
%% answer: BC
%% memo:
\memoanswer{
小球从 D 点运动到 A 点，由动能定理有 $mg(H-2R)=\frac{1}{2}mv_{A}^{2}$，解得 $v_{A}=\sqrt{2g(H-2R)}$，小球运动到 A 点且做平抛运动应满足条件 $v_{A}>0$，故 $H>2R$，选项 C 正确、D 错误；小球经过 A 点做平抛运动，根据平抛运动规律：$x=v_{A}t$，$2R=\frac{1}{2}gt^{2}$，解得 $x=2\sqrt{2RH-4R^{2}}$，选项 A 错误、B 正确。
}

\item
%% number: 19
%% paperName: 2012年浙江理综第 19 题 6 分
%% typeId: 2
%% type: 多选题
%% chapter: 电场
%% point: 静电感应
%% method: 无
%% score: 6
%% degree: 450
%% duplicateId: 0
%% body:
用金属箔做成一个不带电的圆环，放在干燥的绝缘桌面上。小明同学用绝缘材料做的笔套与头发摩擦后，将笔套自上向下慢慢靠近圆环，当距离约为$0.5\Ucm$时圆环被吸引到笔套上，如图所示。对上述现象的判断与分析，下列说法正确的是\xzanswer{ABC}
\onepicture[h=2.4cm, align=right]{figs/19.png}
\fourchoices[answer=ABC]
{摩擦使笔套带电}
{笔套靠近圆环时，圆环上、下部感应出异号电荷}
{圆环被吸引到笔套的过程中，圆环所受静电力的合力大于圆环的重力}
{笔套碰到圆环后，笔套所带的电荷立刻被全部中和}
%% answer: ABC
%% memo:
\memoanswer{
笔套与头发摩擦，笔套带电，故选项 A 正确；带电的笔套靠近金属箔做成的圆环，由于静电感应，在圆环的上、下部感应出异种电荷，故选项 B 正确；圆环被吸引到笔套上的过程中，圆环有向上的加速度，静电力大于圆环的重力，故选项 C 正确；笔套碰到圆环后，笔套上的电荷有一部分转移到圆环，不会立刻被全部中和，故选项 D 错误。
}

\item
%% number: 20
%% paperName: 2012年浙江理综第 20 题 6 分
%% typeId: 2
%% type: 多选题
%% chapter: 交流电
%% point: LC振荡回路
%% method: 无
%% score: 6
%% degree: 450
%% duplicateId: 0
%% body:
为了测量储罐中不导电液体的高度，将与储罐外壳绝缘的两块平行金属板构成的电容器$C$置于储罐中，电容器可通过开关S与线圈$L$或电源相连，如图所示。当开关从a拨到b时，由$L$与$C$构成的回路中产生周期$T=2\pi\sqrt {LC} $的振荡电流。当罐中的液面上升时\xzanswer{BC}
\onepicture[h=3.2cm, align=right]{figs/20.png}
\fourchoices[answer=BC]
{电容器的电容减小}
{电容器的电容增大}
{$LC$回路的振荡频率减小}
{$LC$回路的振荡频率增大}
%% answer: BC
%% memo:
\memoanswer{
由于罐中的液体是不导电的，介电常数比空气大。当液面上升时，金属板间的电介质的介电常数增加，因此，电容器的电容增大，选项 A 错误、B 正确；根据振荡电路的周期公式可知，振荡周期增大，振荡频率减小，选项 D 错误、C 正确。
}

\item
%% number: 21
%% paperName: 2012年浙江理综第 21 题 10 分
%% typeId: 4
%% type: 实验题
%% chapter: 光学
%% point: 测定玻璃折射率
%% method: 无
%% score: 10
%% degree: 650
%% duplicateId: 0
%% body:
在“测定玻璃的折射率”实验中，某同学经正确操作插好了4枚大头针，如图甲所示。
\twopicture[showlabel=false, h=3.0cm, align=right]
{figs/21a.png}
{figs/21b.png}
\begin{enumerate}
	\item
	在图中画出完整的光路图；
	\item
	对你画出的光路图进行测量和计算，求得该玻璃砖的折射率$n=$\_\_\_\_\_\_\_\_（保留3位有效数字）；
	\item
	为了观测光在玻璃不同表面的折射现象，某同学做了两次实验，经正确操作插好了8枚大头针，如图乙所示。图中P$_{1}$和P$_{2}$是同一入射光线上的2枚大头针，其对应出射光线上的2枚大头针是P$_{3}$和\_\_\_\_\_\_\_\_（填“A”或“B”）。
\end{enumerate}
%% answer: 
\jdanswer{
\begin{enumerate}
	\item
	光路图如图所示（见解析）。

	\item
	$1.51$（$1.48\sim1.54$）

	\item
	A

\end{enumerate}
}
%% memo:
\memoanswer{
（1）完整的光路图如图所示。
\onepicture[h=3.0cm, align=right]{figs/21_an.png}
（2）用量角器测出入射角 $i$ 与折射角 $r$，根据折射定律 $n=\frac{\sin i}{\sin r}$ 得出结果；或者利用坐标纸结合入射角和折射角画两个直角三角形，然后用刻度尺测出所对应的直角边和斜边的长度，进一步计算出正弦值，再代入折射定律公式即可。

（3）由图可知，由于入射光线比较靠近玻璃砖的右边，经过上表面折射后，再经过右侧面折射出来，故应该通过 A 点。
}

\item
%% number: 22
%% paperName: 2012年浙江理综第 22 题 10 分
%% typeId: 4
%% type: 实验题
%% chapter: 力物体的平衡
%% point: 验证平行四边形实验
%% method: 无
%% score: 10
%% degree: 650
%% duplicateId: 0
%% body:
在“探究求合力的方法”实验中，现有木板、白纸、图钉、橡皮筋、细绳套和一把弹簧秤。
\begin{enumerate}
	\item
	为完成实验，某同学另找来一根弹簧，先测量其劲度系数，得到的实验数据如下表：

\begin{center}
\begin{tabular}{|c|c|c|c|c|c|c|c|}
\hline
弹力$F$（$\UN$） & 0.50 & 1.00 & 1.50 & 2.00 & 2.50 & 3.00 & 3.50 \\
\hline
伸长量$x$（$10^{-2}\Um$） & 0.74 & 1.80 & 2.80 & 3.72 & 4.60 & 5.58 & 6.42 \\
\hline
\end{tabular}
\end{center}

	用作图法求得该弹簧的劲度系数$k=$\_\_\_\_\_\_\_\_$\UNm$；
	\onepicture[h=3.6cm, align=right]{figs/22a.png}
	\item
	某次实验中，弹簧秤的指针位置如图所示，其读数为\_\_\_\_\_\_\_\_$\UN$；同时利用（1）中结果获得弹簧上的弹力值为$2.50\UN$，请在答题纸上画出这两个共点力的合力$F_{\text{合}}$；
	\onepicture[h=3.0cm, align=right]{figs/22b.png}
	\item
	由图得到$F_{\text{合}}=$\_\_\_\_\_\_\_\_$\UN$。
\end{enumerate}
%% answer: 
\jdanswer{
\begin{enumerate}
	\item
	$53$（$51\sim55$）

	\item
	$2.10$（$2.08\sim2.12$），图略

	\item
	$3.3$（$3.1\sim3.5$）

\end{enumerate}
}
%% memo:
\memoanswer{
（1）根据表格数据描点，然后连成一条过原点的直线，如图所示，直线的斜率等于弹簧的劲度系数，在图象上取相距较远的两点 $P(0.50,0.31)$、$Q(6.00,3.25)$，$k=\frac{3.25-0.31}{6.00-0.50}\UNcm\approx0.53\UNcm=53\UNm$。
\onepicture[h=3.6cm, align=right]{figs/22_an1.png}

（2）读出弹簧测力计的读数为 $2.10\UN$（保留三位有效数字）。以 O 为顶点，画出两弹簧的绳套方向就是两拉力方向，再确定并画好力的标度，画出两拉力的图示，以两拉力为邻边作出平行四边形，画出平行四边形的对角线，即合力 $F_{\text{合}}$ 如图所示。
\onepicture[h=3.2cm, align=right]{figs/22_an2.png}

（3）用刻度尺量出合力的长度，根据确定的标度算出合力的大小。
}

\item
%% number: 23
%% paperName: 2012年浙江理综第 23 题 16 分
%% typeId: 6
%% type: 计算题
%% chapter: 运动和力
%% point: 牛顿定律的应用
%% method: 无
%% score: 16
%% degree: 450
%% duplicateId: 0
%% body:
为了研究鱼所受水的阻力与其形状的关系，小明同学用石蜡做成两条质量均为$m$、形状不同的“A鱼”和“B鱼”，如图所示。在高出水面$H$处分别静止释放“A鱼”和“B鱼”，“A鱼”竖直下潜$h_{A}$后速度减为零，“B鱼”竖直下潜$h_{B}$后速度减为零。“鱼”在水中运动时，除受重力外，还受浮力和水的阻力。已知“鱼”在水中所受浮力是其重力的$\frac{10}{9}$倍，重力加速度为$g$，“鱼”运动的位移值远大于“鱼”的长度。假设“鱼”运动时所受水的阻力恒定，空气阻力不计。求：
\begin{enumerate}
	\item
	“A鱼”入水瞬间的速度$v_{A1}$；
	\item
	“A鱼”在水中运动时所受阻力$f_{A}$；
	\item
	“A鱼”与“B鱼”在水中运动时所受阻力之比$f_{A}:f_{B}$。
\end{enumerate}
\onepicture[h=4.0cm, align=right]{\includesvg{figs/23}}
%% answer: 
\jdanswer{
\begin{enumerate}
	\item
	$\sqrt {2gH} $

	\item
	$mg$（$\frac{H}{{{h_A}}} - \frac{1}{9}$）

	\item
	$\frac{{{h_B}(9H - {h_A})}}{{{h_A}(9H - {h_B})}}$

\end{enumerate}
}
%% memo:
\memoanswer{
（1）“A 鱼”在入水前做自由落体运动，有 $v_{A1}^{2}-0=2gH$，得 $v_{A1}=\sqrt{2gH}$。

（2）“A 鱼”在水中运动时受重力、浮力和阻力的作用，做匀减速运动，设加速度为 $a_{A}$，有 $F_{\text{合}}=F_{\text{浮}}+f_{A}-mg$，$F_{\text{合}}=ma_{A}$，$0-v_{A1}^{2}=-2a_{A}h_{A}$，由题意 $F_{\text{浮}}=\frac{10}{9}mg$。综合上述各式，得 $f_{A}=mg\left(\frac{H}{h_{A}}-\frac{1}{9}\right)$。

（3）考虑到“B 鱼”的受力、运动情况与“A 鱼”相似，有 $f_{B}=mg\left(\frac{H}{h_{B}}-\frac{1}{9}\right)$，解得 $\frac{f_{A}}{f_{B}}=\frac{h_{B}(9H-h_{A})}{h_{A}(9H-h_{B})}$。
}

\item
%% number: 24
%% paperName: 2012年浙江理综第 24 题 20 分
%% typeId: 6
%% type: 计算题
%% chapter: 磁场
%% point: 带电粒子在复合场中的运动
%% method: 无
%% score: 20
%% degree: 450
%% duplicateId: 0
%% body:
如图所示，两块水平放置、相距为$d$的长金属板接在电压可调的电源上。两板之间的右侧区域存在方向垂直纸面向里的匀强磁场。将喷墨打印机的喷口靠近上板下表面，从喷口连续不断喷出质量均为$m$、水平速度均为$v_{0}$、带相等电荷量的墨滴。调节电源电压至$U$，墨滴在电场区域恰能沿水平向右做匀速直线运动；进入电场、磁场共存区域后，最终垂直打在下板的$M$点。
\begin{enumerate}
	\item
	判断墨滴所带电荷的种类，并求其电荷量；
	\item
	求磁感应强度$B$的值；
	\item
	现保持喷口方向不变，使其竖直下移到两板中间的位置。为了使墨滴仍能到达下板$M$点，应将磁感应强度调至$B'$，则$B'$的大小为多少？
\end{enumerate}
\onepicture[h=2.4cm, align=right]{figs/24.png}
%% answer: 
\jdanswer{
\begin{enumerate}
	\item
	$q=\frac{mgd}{U}$，墨滴带负电荷

	\item
	$B=\frac{v_{0}U}{gd^{2}}$

	\item
	$B'=\frac{4v_{0}U}{5gd^{2}}$

\end{enumerate}
}
%% memo:
\memoanswer{
（1）墨滴在电场区域做匀速直线运动，有 $q\frac{U}{d}=mg$，解得 $q=\frac{mgd}{U}$。由于电场方向向下，电荷所受电场力向上，可知墨滴带负电荷。

（2）墨滴垂直进入电、磁场共存区域，重力仍与电场力平衡，合力等于洛伦兹力，墨滴做匀速圆周运动，有 $qv_{0}B=m\frac{v_{0}^{2}}{R}$。考虑墨滴进入磁场和撞板的几何关系，可知墨滴在该区域恰完成四分之一圆周运动，则半径 $R=d$。联立解得 $B=\frac{v_{0}U}{gd^{2}}$。

（3）根据题设，墨滴运动轨迹如图，设圆周运动半径为 $R'$，有 $qv_{0}B'=m\frac{v_{0}^{2}}{R'}$。由图示可得 $R'^{2}=d^{2}+\left(R'-\frac{d}{2}\right)^{2}$，得 $R'=\frac{5}{4}d$。
\onepicture[h=3.0cm, align=right]{figs/24_an.png}

联立解得 $B'=\frac{4v_{0}U}{5gd^{2}}$。
}

\item
%% number: 25
%% paperName: 2012年浙江理综第 25 题 22 分
%% typeId: 6
%% type: 计算题
%% chapter: 电磁感应
%% point: 切割磁感线电动势
%% method: 无
%% score: 22
%% degree: 250
%% duplicateId: 0
%% body:
为了提高自行车夜间行驶的安全性，小明同学设计了一种“闪烁”装置。如图所示，自行车后轮由半径$r_{1}=5.0\times10^{-2}\Um$的金属内圈、半径$r_{2}=0.40\Um$的金属外圈和绝缘辐条构成。后轮的内、外圈之间等间隔地接有4根金属条，每根金属条的中间均串联有一电阻值为$R$的小灯泡。在支架上装有磁铁，形成了磁感应强度$B=0.10\UT$、方向垂直纸面向外的“扇形”匀强磁场，其内半径为$r_{1}$、外半径为$r_{2}$、张角$\theta=\frac{\pi }{6}$。后轮以角速度$\omega=2\pi\Urads$相对于转轴转动。若不计其他电阻，忽略磁场的边缘效应。
\begin{enumerate}
	\item
	当金属条$ab$进入“扇形”磁场时，求感应电动势$E$，并指出$ab$上的电流方向；
	\item
	当金属条$ab$进入“扇形”磁场时，画出“闪烁”装置的电路图；
	\item
	从金属条$ab$进入“扇形”磁场时开始，经计算画出轮子转一圈过程中，内圈与外圈之间电势差$U_{ab}$随时间$t$变化的$U_{ab}-t$图象；
	\item
	若选择的是“1.5V、0.3A”的小灯泡，该“闪烁”装置能否正常工作？有同学提出，通过改变磁感应强度$B$、后轮外圈半径$r_{2}$、角速度$\omega$和张角$\theta$等物理量的大小，优化前同学的设计方案，请给出你的评价。
\end{enumerate}
\onepicture[h=3.4cm, align=right]{figs/25.png}
%% answer: 
\jdanswer{
\begin{enumerate}
	\item
	$E=4.9\times10^{-2}\UV$，电流方向为$b\to a$

	\item
	电路图如图所示（见解析）

	\item
	$U_{ab}-t$图象如图所示（见解析）

	\item
	不能正常工作；$B$、$r_{2}$、$\omega$增大均使$E$增大但有限度，$\theta$增大时$E$不变。

\end{enumerate}
}
%% memo:
\memoanswer{
（1）金属条 ab 在磁场中切割磁感线时，所构成的回路的磁通量变化。设经过时间 $\Delta t$，磁通量变化量为 $\Delta\Phi$，由法拉第电磁感应定律 $E=\frac{\Delta\Phi}{\Delta t}$，$\Delta\Phi=B\Delta S=B\left(\frac{1}{2}r_{2}^{2}\Delta\theta-\frac{1}{2}r_{1}^{2}\Delta\theta\right)$，联立解得 $E=\frac{\Delta\Phi}{\Delta t}=\frac{1}{2}B\omega(r_{2}^{2}-r_{1}^{2})=4.9\times10^{-2}\UV$。根据右手定则（或楞次定律），可得感应电流方向为 $b\to a$。

（2）通过分析，可得电路图为
\onepicture[h=2.6cm, align=right]{figs/25_an1.png}

（3）设电路中的总电阻为 $R_{\text{总}}$，根据电路图可知 $R_{\text{总}}=R+\frac{1}{3}R=\frac{4}{3}R$。ab 两端电势差 $U_{ab}=E-IR=E-\frac{E}{R_{\text{总}}}R=\frac{1}{4}E=1.2\times10^{-2}\UV$。设 ab 离开磁场区域的时刻为 $t_{1}$，下一根金属条进入磁场区域的时刻为 $t_{2}$，则 $t_{1}=\frac{\theta}{\omega}=\frac{1}{12}\Us$，$t_{2}=\frac{\pi}{\omega}=\frac{1}{4}\Us$。设轮子转一圈的时间为 $T$，则 $T=\frac{2\pi}{\omega}=1\Us$。在 $T=1\Us$ 内，金属条有四次进出，后三次与第一次相同。可画出 $U_{ab}-t$ 图象如图所示。
\onepicture[h=2.6cm, align=right]{figs/25_an2.png}

（4）“闪烁”装置不能正常工作。（金属条的感应电动势只有 $4.9\times10^{-2}\UV$，远小于小灯泡的额定电压，因此无法工作。）

$B$ 增大，$E$ 增大，但有限度；

$r_{2}$ 增大，$E$ 增大，但有限度；

$\omega$ 增大，$E$ 增大，但有限度；

$\theta$ 增大，$E$ 不变。
}

\item
%% number: 33
%% paperName: 2012年浙江理综第 33 题 10 分
%% typeId: 3
%% type: 填空题
%% chapter: 气体性质
%% point: 气体图象
%% method: 无
%% score: 10
%% degree: 450
%% duplicateId: 0
%% body:
一定质量的理想气体，状态从A$\to$B$\to$C$\to$D$\to$A的变化过程可用如图所示的$p-V$的图描述，图中$p_{1}$、$p_{2}$、$V_{1}$、$V_{2}$和$V_{3}$为已知量。
\begin{enumerate}
	\item
	气体状态从A到B是\_\_\_\_\_\_过程（填“等容”、“等压”或“等温”）；
	\item
	状态从B到C的变化过程中，气体的温度\_\_\_\_\_\_（填“升高”、“不变”或“降低”）；
	\item
	状态从C到D的变化过程中，气体\_\_\_\_\_\_（填“吸热”或“放热”）；
	\item
	状态从A$\to$B$\to$C$\to$D的变化过程中，气体对外界所做的总功为\_\_\_\_\_\_\_\_。
\end{enumerate}
\onepicture[h=3.4cm, align=right]{figs/33.png}
%% answer: 
\jdanswer{
\begin{enumerate}
	\item
	等压

	\item
	降低

	\item
	放热

	\item
	$p_{1}$（$V_{3}-V_{1}$）$-p_{2}$（$V_{3}-V_{2}$）

\end{enumerate}
}
%% memo:
\memoanswer{
（1）由题图可得从 A 到 B 压强不变；（2）从 B 到 C 的变化过程中体积不变，压强减小，由 $p/T=C$，温度降低；（3）从 C 到 D 的变化过程中压强不变，体积减小，外界对气体做功，而 $V/T=C$，温度降低，由 $\Delta U=W+Q$，气体放热。（4）由于 $W=F\Delta l=pS\Delta l=p\Delta V$，从 A 到 B 气体体积增大对外界做功 $W_{1}=p_{1}(V_{3}-V_{1})$；从 B 到 C 气体体积不做功；从 C 到 D 体积减小，外界对气体做功 $W_{2}=p_{2}(V_{3}-V_{2})$。从 $A\to B\to C\to D$ 的变化过程中，气体对外界所做的总功为 $W=W_{1}-W_{2}=p_{1}(V_{3}-V_{1})-p_{2}(V_{3}-V_{2})$。
}

\item
%% number: 34
%% paperName: 2012年浙江理综第 34 题 10 分
%% typeId: 3
%% type: 填空题
%% chapter: 原子和原子核
%% point: 质能方程
%% method: 无
%% score: 10
%% degree: 650
%% duplicateId: 0
%% body:
一个静止的钚核\ce{^239_94Pu}自发衰变成一个铀核\ce{^235_92U}和另一个原子核X，并释放出一定的能量。其核衰变方程为：\ce{^239_94Pu -> ^235_92U + X}。
\begin{enumerate}
	\item
	方程中的“X”核符号为\_\_\_\_\_\_；
	\item
	钚核的质量为$239.0522\Uu$，铀核的质量为$235.0439\Uu$，X核的质量为$4.0026\Uu$，已知$1\Uu$相当于$931\UMeV$，则该衰变过程放出的能量是\_\_\_\_\_\_$\UMeV$；
	\item
	假设钚核衰变释放出的能量全部转变为铀核和X核的动能，则X核与铀核的动能之比是\_\_\_\_\_\_。
\end{enumerate}
%% answer: 
\jdanswer{
\begin{enumerate}
	\item
	\ce{^4_2He}

	\item
	$5.31$（$\pm0.02$ 都可）

	\item
	$58.7$（$\pm0.1$ 都可）

\end{enumerate}
}
%% memo:
\memoanswer{
（1）由质量守恒和电荷量守恒可得：X 为氦原子核 \ce{^4_2He}。

（2）由 $\Delta E=\Delta mc^{2}=(239.0522-235.0439-4.0026)\times931\UMeV=5.31\UMeV$。

（3）由于自发衰变动量守恒，则有 $m_{U}v_{U}+m_{X}v_{X}=0$，钚核衰变释放出的能量全部转变为铀核和 X 核的动能，$\frac{1}{2}m_{U}v_{U}^{2}+\frac{1}{2}m_{X}v_{X}^{2}=\Delta E$，X 核与铀核的动能之比 $\frac{\frac{1}{2}m_{X}v_{X}^{2}}{\frac{1}{2}m_{U}v_{U}^{2}}=58.7$。
}

\end{enumerate}

\end{document}
